\section{Method: Entropy-based Router Consolidation}
\label{sec:method}

\subsection{SR-Core Architecture}
\label{sec:arch}

We briefly formalize the architecture to fix notation.
A block bank $\Bcal = \{F_1, \ldots, F_n\}$ contains $n$ parameter blocks, each a two-layer
MLP with shared input/output dimension $d$.
The model processes each input token through $R$ recursive applications of the bank.
At step $r = 1$, a learned router produces a probability distribution over blocks:
\begin{equation}
  \mathbf{p}_t
    = \mathrm{softmax}\!\left(\mathbf{q}_t^\top \mathbf{K}\right) \in \mathbb{R}^n
  \label{eq:router}
\end{equation}
where $\mathbf{q}_t = W_q \mathbf{h}_{t,0}$ is a query derived from the initial hidden
state and $\mathbf{K} \in \mathbb{R}^{n \times d_k}$ is a learned key matrix.
The \textbf{SR-Core} selection is
\begin{equation}
  S_t = \topk(\mathbf{p}_t), \qquad \WS = |S_t| = k
  \label{eq:srcore}
\end{equation}
with soft gates $g_{t,b} \propto p_{t,b}$ for $b \in S_t$.
The hidden state update at step $r$ is
\begin{equation}
  \mathbf{h}_{t,r}
    = \mathbf{h}_{t,r-1}
      + \sum_{b \in S_t} g_{t,b} \cdot F_b(\mathbf{h}_{t,r-1})
  \label{eq:update}
\end{equation}
using the \emph{same} selection $S_t$ for all $r = 1, \ldots, R$.
The active working set is therefore exactly $k$ blocks per token, independent of $R$ and $n$.

Training uses a deep-supervision language modeling loss that supervises the output at every
recursion depth, with end-weighted contributions to preserve full-depth quality while
allowing intermediate readouts.

\subsection{Entropy-based Consolidation Objective}
\label{sec:objective}

We augment the language modeling loss with an entropy penalty on the pre-selection router
distribution:
\begin{equation}
  \Lcal = \Lcal_{\mathrm{LM}} + \lambda \cdot \Hcal
  \label{eq:loss}
\end{equation}
where
\begin{equation}
  \Hcal = \frac{1}{N} \sum_{t=1}^{N} H\!\left(\mathbf{p}_t^{(1)}\right),
  \qquad
  H(\mathbf{p}) = -\sum_{i=1}^{n} p_i \log p_i
  \label{eq:entropy}
\end{equation}
and $\mathbf{p}_t^{(1)}$ is the router distribution for token $t$ at $r = 1$.
The entropy penalty $\Hcal$ encourages the router to assign concentrated probability mass
before the top-$k$ selection: lower entropy corresponds to higher confidence in the selected
blocks, which in practice produces more consistent routing choices across tokens with similar
content and thereby improves cross-token cache hit rates.

\paragraph{Why $r = 1$ only.}
Routing decisions are made at $r = 1$; steps $r = 2, \ldots, R$ reuse the same block set
$S_t$ and do not have a meaningful ``unconcentrated'' distribution to regularize.
Applying the entropy loss at $r = 1$ is therefore the canonical intervention point.

\paragraph{Why pre-top-$k$.}
The router distribution $\mathbf{p}_t$ is computed before the top-$k$ cut.
Applying entropy regularization at this point shapes the entire selection landscape: the
penalty encourages the router to place probability mass on a small subset of blocks before
any discretization, rather than penalizing the post-selection gates (which are already
sparse by construction).
A concentrated pre-selection distribution also improves the router's top-$k$ margin
$p_k - p_{k+1}$, stabilizing which blocks fall at the selection boundary.

\paragraph{Distinction from pairwise objectives.}
We test soft-Jaccard-style overlap objectives as an internal negative control.
Such objectives operate on \emph{cross-token} routing similarity: they compare the block
selections of consecutive tokens and penalize dissimilarity.
This is a different quantity from the \emph{within-token} distribution concentration that
entropy targets.
In our evaluation, soft-Jaccard objectives increase soft overlap metrics while
simultaneously increasing the number of unique routing combinations --- the opposite of the
cache-beneficial consolidation we seek, and the wrong lever for controlling pre-selection
sharpness.
Entropy minimization, by acting on the per-token distribution before the top-$k$ cut,
produces the qualitatively different effect of a reduced unique routing vocabulary.

\paragraph{Relationship to load balancing.}
Standard MoE load-balancing losses \citep{fedus2022switch} promote \emph{uniform} block
utilization, which is the opposite of entropy minimization.
We do not replace load balancing with the entropy objective; the two terms address different
failure modes (dead blocks vs.\ over-diffuse routing) and can be combined.
In our experiments we observe that trained SR-Core models already activate all $n$ blocks
across a corpus; load balancing is not required as a prerequisite.

\paragraph{Hyperparameter $\lambda$.}
We sweep $\lambda \in \{0.001, 0.003, 0.005, 0.007\}$.
The penalty is dimensionally compatible with the language modeling loss (both in nats),
making $\lambda$ a direct ratio between the two terms.
Across our sweep, $\lambda = 0.003$ is the most reliable generalization-balanced operating
point, with cache and code-generalization improvements that replicate across seeds.
$\lambda = 0.005$ is the stronger cache/systems point, delivering larger LRU byte
reductions; its quality effects are seed-dependent and should not be characterized as a
uniform generalization cost (Section~\ref{sec:results}).
